% Research proof fragment, 7 October 2026. Not an assertion of priority.
% This file is intentionally separate from the main manuscript until reviewed.
\section{Retaining the old support in the bad third-prefix boxes}
\label{sec:retained-base-upper}

Set $\theta_0=89/160$ and $\theta=29/50$. Let $\Phi$ be the full
eight-branch envelope of Lichtman 2023, Proposition 6.1, (6.4), at
$\alpha=1$, $\sigma=7/32$, as in Section~\ref{sec:sharper-upper-support}.
Write $S=t_1+t_2+t_3$ and
\[
 E_\theta=\{0<t_3\le t_2\le t_1<1/5:\quad
 3t_1\le\theta,\quad t_1+t_2+3t_3\le\theta,\quad
 \Phi(t_1,t_2,t_3)<\theta\}.
\]
Extend the ordinary linear lower-sieve function by $f(s)=0$ for
$s\le2$. Define
\begin{equation}\label{eq:retained-base-functional}
 K(\theta)=\frac4\theta+2e^{-\gamma}
 \int_{E_\theta}
 \frac{f((\theta-S)/t_3)-f((\theta_0-S)/t_3)}
 {t_1t_2t_3^2}\,dt_1dt_2dt_3.
\end{equation}
No numerical assertion is part of this definition.

\begin{proposition}[Retained-base upper-sieve interface]
\label{prop:retained-base-upper}
For every $K_*>K(29/50)$, the conclusions of
Proposition~\ref{prop:pruned-upper} hold with $K_*$ in place of
$3549/500$. The product families, multiplicities, moving residue $N$,
and full fourth-prime range are exactly those in that proposition.
\end{proposition}

We first justify the asymptotic equality needed for the residual lower
sieves. A one-sided lower-sieve bound alone would not justify subtracting
two residual density sums.

\begin{lemma}[Uniform residual lower density]
\label{lem:uniform-residual-lower-density}
Let $R=N^\beta$, $y=N^t$, where $(\beta,t)$ ranges over a compact
subset of $\{\beta>0,t>0,\beta\ge2t\}$. Put
$g_N(p)=\mathbf1_{p>2,\,p\nmid N}/(p-1)$ and
$V_N(y)=\prod_{p<y}(1-g_N(p))$. With the ordinary Rosser support
\[
 \mathcal D^-(R)=\{d=q_1\cdots q_r:\ q_1>\cdots>q_r,
 q_1^2<R,\quad q_1\cdots q_{j-1}q_j^3<R
 \text{ for even }j\le r\},
\]
including the empty product, one has uniformly
\[
 \frac{\sum_{d\mid P_N(y),\ d\in\mathcal D^-(R)}\mu(d)g_N(d)}
 {V_N(y)}=f(\beta/t)+o(1).
\]
The strict or non-strict cube convention gives the same asymptotic.
\end{lemma}
\begin{proof}
Here is the precise classical input used, including convergence. For
dimension one and the cubic Rosser support, the normalized positive
first-exit integrals $a_j(s)$ obey
\[
 a_1(s)=(3-s)_+/s,\qquad
 a_j(s)=\frac1s\int_{\max(s,\,2+\epsilon_j)}^{2+j}
                  a_{j-1}(u-1)\,du\quad(j\ge2),
\]
where $\epsilon_j=0$ for even $j$ and $1$ for odd $j$, and an
empty integral is zero. The classical Rosser calculation gives
\begin{equation}\label{eq:classical-positive-exit-series}
 \sum_{j\ge2,\ j\ {
m even}}a_j(s)=1-f(s)\quad(s\ge2),
\end{equation}
with uniform convergence on compact $s$ intervals. This is a
convergence input, not a conclusion drawn merely from finite-depth
asymptotics. Explicit sources are
\href{https://www.suzuyu1729.com/notes/NagoyaSeminar2022/RosserIwaniec.pdf}
{Y. Suzuki, \emph{The Rosser--Iwaniec sieve}}, Sections 7--9
(the exit recurrence and (9.1)--(9.2)), Propositions 9.4 and 11.8
(compact convergence and its existence), and Section 12
(the optimal linear parameter $\beta=2$).
The normalization to $f$ is also given by (3.2.10) in
\href{https://mathweb.tifr.res.in/Documents/Publications/Lectures/tifr72.pdf}
{Y. Motohashi, \emph{Lectures on Sieve Methods and Prime Number Theory}},
TIFR 72 (1983); its Section 3.3, Lemma 14, supplies the convergence
estimate in the classical linear case. These inputs concern the
ordinary support displayed here, not a factorable modified weight.

Since $y\le\sqrt R$, the initial square condition never excludes
a prime below $y$. All first exclusions therefore have even length.
The finite first-exit identity gives exactly
\begin{equation}\label{eq:residual-lower-positive-exits}
 \frac{V_N^-(R,y)}{V_N(y)}
 =1-\sum_{j\ge2,\ j\ {
m even}} T_{j,N}(R,y),\qquad
 T_{j,N}=\sum_{\substack{P=q_1\cdots q_j\text{ first exit}\\q_1<y}}
 g_N(P)\frac{V_N(q_j)}{V_N(y)}\ge0.
\end{equation}
Fix a rank $j$. Every preceding prefix belongs to
$\mathcal D^{\rm well}(R)$: the even cube condition implies its
own square condition and, by ordering, the next square condition;
the first square is part of the lower support definition.
If $A_i=q_1\cdots q_i$, then $A_i^2<R A_{i-1}$, and induction gives
$A_i<R^{1-2^{-i}}$. Failure of the $j$th cube consequently forces
\[
 q_j^3\ge R/A_{j-1}>R^{2^{1-j}}.
\]
Thus all primes in this fixed-depth sum exceed $N^{c_j}$ for some
$c_j>0$ uniform in the compact parameter set.

For $g_0(p)=\mathbf1_{p>2}/(p-1)$, prime-number-theorem integration
on this compact exponent domain gives
$T_{j,0}(R,y)=a_j(\beta/t)+o(1)$ uniformly. The recurrence above
can alternatively be obtained by first summing the largest prime.
All boundary hyperplanes have zero limiting measure.
Deleting $p\mid N$ changes this finite-depth asymptotic by $o(1)$:
\[
 \sum_{p\mid N,\ p\ge N^{c_j}}\frac1{p-1}
 \ll N^{-c_j}\log N,
\]
and the other prime reciprocal sums and the product ratio are bounded
in terms of $j$ and the compact parameter set. More explicitly, each
surviving normalized summand satisfies
\[
 g_N(P)\frac{V_N(q_j)}{V_N(y)}
 =g_0(P)\frac{V_0(q_j)}{V_0(y)}
   \prod_{\substack{q_j\le p<y\\p\mid N}}(1-g_0(p))
 =\bigl(1+o(1)\bigr)
   g_0(P)\frac{V_0(q_j)}{V_0(y)}.
\]
Consequently $T_{j,N}=a_j(\beta/t)+o(1)$ for each fixed $j$.

There is no interchange of an uncontrolled infinite prime sum here.
The standard one-sided linear lower-density bound applies uniformly
to $g_N$, since deleting primes decreases every inverse local-product
ratio. It gives $V_N^-/V_N\ge f(\beta/t)-o(1)$.
In the opposite direction, discard the positive exits of rank $>J$
from \eqref{eq:residual-lower-positive-exits}. The fixed-depth
asymptotics then give
\[
 \frac{V_N^-}{V_N}\le
 1-\sum_{2\le j\le J,\ j\ {
m even}}a_j(\beta/t)+o_J(1).
\]
First let $N\to\infty$, then $J\to\infty$, using the compact-uniform
convergence in \eqref{eq:classical-positive-exit-series}. This proves
the lemma, including $\beta/t=2$ by continuity. In particular no
two-sided dimension hypothesis with a constant independent of $N$
has been silently assumed.
\end{proof}

\begin{proof}[Proof of Proposition~\ref{prop:retained-base-upper}]
\paragraph{Levels and the full envelope.}
Fix $e>0$ sufficiently small and use the two raw exponents
\[
 \nu_0=\theta_0-2e,\qquad \nu_1=\theta-2e,
 \qquad R_i=N^{\nu_i},\qquad \theta_A=\theta-e.
\]
Take the small-prime and rough windows, beta weights, exponent grid,
rank cutoff and hierarchy of Section~\ref{sec:pruned-upper-sieve},
with $D_0=R_1$. Thus $w=N^\omega$, the beta level is $N^\ell$,
the grid mesh is $\xi\ll\omega^5$, and every parameter is fixed
before $N\to\infty$.

Here is a buffered box rule which uses all eight branches and handles
their gates. Set $v(t)$ as in Section~\ref{sec:sharper-upper-support},
and write $W_1,W_2$ for the inner minima in
$w(t_1,t_2,t_3)$ and $w(t_1,t_3,t_2)$, respectively.
A closed ordered third-prefix box is good if throughout it at least
one fixed member of the following finite list supplies its indicated
test:
\[
 \begin{gathered}
 v(t_1),\ v(t_2),\ v(t_1+t_2),\ v(S)\ \ge\theta_A;\\
 \max\{v(t_1),W_i\}\ge\theta_A\quad(i=1,2);\\
 v(t_1+t_3)\ge\theta_A,\quad
 v(t_1+t_3)-(t_1+2t_2+t_3)\ge e;\\
 v(t_2+t_3)\ge\theta_A,\quad
 v(t_2+t_3)-(2t_1+t_2+t_3)\ge e.
 \end{gathered}
\]
Otherwise the box is bad. The label decision is made once and is
independent of the eventual rank and its parity. Third primes below
$w$ have the good label. For small enough $\omega,\xi$, this agrees
with the continuous $W_1$ test on all relevant cube-supported vectors.

The first raw cube gives $t_1\le\theta/3+O(\xi)<25/128$ with a
fixed gap. Hence $v(t_1)=(1+t_1)/2$, which is exactly the gate of
both $W_i$. The maximum with that gate removes the discontinuity
of each $w$ branch. The pair branches instead have the explicit
positive gate margin $e$; their discontinuities are not treated as
continuous. Scale changes and changing the residue exponent to
$\alpha=1+3\kappa$ perturb their affine pieces by $O(\kappa)$.
They therefore retain an envelope value above
$\theta-3e/2$ when $\kappa\ll e$. If desired, for each pair atom
one can choose its factorization parameter a little below its
$v$-value and above its gate, rather than use a common parameter
below the gate. Both choices have surplus over its actual raw
support. At a $W_i$ gate crossing the $v(t_1)$ branch supplies the
same surplus. This explicitly preserves the hypotheses in the proof
of Lichtman 2023, Proposition 6.1.

Bad boxes satisfy, up to $O(\xi)$,
\[
 t_1\le2\theta_A-1,\qquad
 t_3\ge(153/32-8\theta_A)/3.
\]
For the latter assertion, on the entire cube-supported domain
$1-2t_2>153/256$ and
$(66-34t_3)/107>153/256-3t_3/8$, so $W_1$ is its middle
affine piece. In particular, in a bad box
\[
 3t_1\le3(2\theta_A-1)+O(\xi)<\nu_0.
\]
This strict old first-cube margin is important below.

\paragraph{The actual mixed support and its upper property.}
Let $\mathcal R$ be the union of $\mathcal D^+(R_0)$ and
the subset of $\mathcal D^+(R_1)$ whose third label is good,
with ranks below three always kept at the new level. Equivalently,
good descendants use $R_1$ and bad descendants use $R_0$.
The support is prefix-closed. A first exclusion is odd: an ordinary
new odd cube fails before the third choice, or a bad third prefix
fails the old third cube, or its descendant first fails a later old
odd cube. An old first-cube failure cannot first appear at a bad
third prefix because of the strict margin just proved.
Thus the first-exit identity \eqref{eq:odd-first-exit-identity}
holds for $\lambda_{\mathcal R}(d)=\mu(d)\mathbf1_{d\in\mathcal R}$.
It proves the pointwise upper majorant and the rough/full one-sided
density comparison \eqref{eq:rough-density-one-sided} for this support.

The exact full density identity is
\begin{equation}\label{eq:retained-base-exact-density}
 V_{\mathcal R}-V_N^+(R_1,z)
 =\sum_{\substack{P=p_1p_2p_3\in\mathcal D^+(R_1)\\
        p_1>p_2>p_3,\ p_1<z\\\text{third label bad}}}
 g_N(P)\left\{V_N^-(R_1/P,p_3)
 -\mathbf1_{P\in\mathcal D^+(R_0)}V_N^-(R_0/P,p_3)\right\}.
\end{equation}
Indeed, the two supports differ only in the disjoint trees with a
bad third prefix. For either raw level $R_i$ at which that prefix is
present, its total descendant contribution is
$-g_N(P)V_N^-(R_i/P,p_3)$: three prefix primes reverse the
parity, so the later odd cubes become the residual even cubes.
The residual initial-square condition is automatic, since
$Pp_3^2<R_i$. This proves the identity with its stated sign.
The old prefix is absent precisely when its third cube fails,
equivalently when $\log(R_0/P)/\log p_3\le2$ (apart from harmless
endpoint conventions). The initial old cube has already been checked.

\paragraph{Boxing, beta weights, and distribution.}
Use exactly the parity boxing in Section~\ref{sec:pruned-upper-sieve},
but with the effective exponent $\nu_1$ on good third labels and
$\nu_0$ on bad labels. Ranks below three use $\nu_1$.
At even ranks impose all relevant lower-endpoint odd cubes with
non-strict inequality; at odd ranks impose the upper-endpoint
odd cubes with strict inequality. Hence $c(d)\ge\lambda_{\mathcal R}(d)$
coefficientwise. Any discrepancy lies within $O(M\xi)$ of an odd
cube boundary at one of the two levels. The same tagged-prime
argument bounds its normalized density cost by $O(\xi\omega^{-5})$.
There is no additional third-label discrepancy.

Append $\beta^+$ to even ranks and $\beta^-$ to odd ranks, with
the outer sign retained. The unique small/rough decomposition gives
$|\rho(q)|\le1$ and the complete pointwise upper bound, exactly as
in \eqref{eq:rough-box-pointwise} and the following beta construction.
Its additional normalized density error is
$O(e^{-c/\eta}\omega^{-2})$.

An actual good atom belongs to
$\mathcal D^{\rm well}(N^{\nu_1+O(M\xi)})$ and has the full
buffered envelope surplus proved above. An actual bad atom belongs
to $\mathcal D^{\rm well}(N^{\nu_0+O(M\xi)})$ and retains its
old first and third cube conditions. Apply
Lemma~\ref{lem:upper-support-89} and the margin argument of
Proposition~\ref{prop:upper-89} to the latter. Choose separate
intermediate factorization parameters for the old and new atoms,
then choose
\[
 \ell+O(M\xi),\ \kappa\ll\epsilon_{\rm disp}\ll\zeta\ll e
\]
small enough for both. Keep one representative slot partition on
each box; the rounded support and appended beta factor cost
$N^{O(M\xi)+\ell}$. This is smaller than the minimum Type II
margin $9\zeta/16$. The reciprocal repeated-label factorial occurs
once, and normalizing each grouped convolution by its slot factorial
gives the bounded coefficients used by the earlier atom argument.
Ranks at most two have complete support below $N^{1/2-c}$.
All other complete moduli have support below $N^{3/5-c}$, since
$\nu_1<.58$. Thus the prime-product distribution proof, its
medium-variable and at-most-two-long-variable alternatives,
squarefree collision estimate, divisor Mellin tails, and uniformity
for residue $N$ apply without alteration. The number of atoms is
fixed before the large-$N$ limit.

In particular the exact coprime mass is replaced by the total mass
with the same bound
$O(X_{\rm out}N^{-a}\log N)$, using $|\rho(q)|\le1$.
The nonnegative outer approximation is global, its excess is counted
once, and the labels and full fourth-prime range are retained.
No distribution estimate for arbitrary bounded outer arithmetic
weights, and no assertion that the whole new weight is triply
well-factorable, is used.

\paragraph{Evaluation and order of limits.}
On all bad boxes, $t_3$ is bounded below by a fixed positive constant:
at the limiting level it is at least $113/2400$. Thus the residual
parameters in \eqref{eq:retained-base-exact-density} range over
compact positive sets wherever a residual term is present. Apply
Lemma~\ref{lem:uniform-residual-lower-density} to both residual sums.
Across the old third-cube boundary its omitted contribution agrees
continuously with $f(s)=0$ at $s=2$.
Uniform Mertens estimates for $V_N(N^{t_3})$ and prime-number-theorem
integration then evaluate the correction in
\eqref{eq:retained-base-exact-density}. Missing prime divisors of
$N$ in these three fixed positive-power prime slots have total
reciprocal mass $o(1)$ uniformly.

The bad sets of the buffered box rules converge in measure to
$E_\theta$. All branches are piecewise affine, and every additional
level or gate equality is contained in a finite union of proper
hyperplanes. Away from these planes, a strict good branch is
eventually certified on its whole box. The $W_i$ rescue and the
explicit pair gate margins above justify this statement despite
the jumps in $\psi$. The correction integrand is bounded on the
compact bad domain, so the extra boundary strips have vanishing cost.

Finally, the ordinary upper density is at most
$V_N(z)(F(5\nu_1)+o(1))$, and $5\nu_1<3$.
Combining all the estimates gives the coefficient
\eqref{eq:retained-base-functional}. First take $N\to\infty$ at
fixed construction parameters; then take the beta and mesh errors
to zero within the analytic margins; finally let $e\downarrow0$.
Given the strict slack $K_*>K(\theta)$, one can choose fixed
parameters whose complete coefficient is below $K_*$.
Applying the resulting upper bound to the five switched masses
$J_i=U_i/8$ and to the bad-semiprime mass $\log(1/b)$ proves
the proposition. This establishes an analytic interface for a
separately certified value of the integral; it makes no claim of
priority or of comparison with every previously published method.
\end{proof}
